\section{Compositional bounded reflection}\label{sec:reflection}

\subsection{The checked one-parameter backend}

The pinned library already contains the essential object proofs for a body \(f(i,x)\):

\begin{itemize}
\item a least-zero search \(\mu_f(N,x)\) with its bound, success, and minimality specification;
\item a universal bit
\[
 \mathsf{allB}_f(N,x)=\prod_{i<N}\operatorname{sg}(f(i,x));
\]
\item an existential bit \(\mathsf{exB}_f(N,x)\) and least nonzero witness;
\item implications between bit equations and object-level bounded quantifiers;
\item decidability of the bounded universal and bounded choice for the existential.
\end{itemize}

The key point is that these declarations already conclude \(\Prf\,\PRAAx\) judgments.  A new
framework should import them, not restate them as future lemmas.

\subsection{Arity-polymorphic search}\label{subsec:searchk}

The present API packages all non-search parameters into one term \(x\).  This is adequate after
pairing, but it creates avoidable encoding overhead and complicates theorem-specialized formulas.
A direct generalization to \(k\) parameters is straightforward and is the first proposed upstream
extension.

Let \(f:\PTm(k+1)\), with variable \(0\) the candidate index and variables
\(1,\ldots,k\) the parameters.  In the recursion step there are \(k+2\) variables:
\(0=N\), \(1=\mu_N\), and \(2,\ldots,k+1\) the parameters.  Define the renaming
\[
 \rho(0)=0,\qquad \rho(j+1)=j+2,
\]
and \(f^{\uparrow}=f[\rho]\).  Put
\[
 s_f(N,m,\vec x)=m+[m=N]\,\operatorname{sg}(f(N,\vec x)).
\]
Primitive recursion with base zero and step \(s_f\) defines
\(\mu_f(N,\vec x)\).

The intended behavior is unchanged: if no zero has appeared before \(N\), then \(m=N\), and the
step advances to \(N+1\) precisely when \(f(N,\vec x)\ne0\).  If a zero was found earlier, then
\(m<N\), the equality bit vanishes, and the value remains fixed.

\begin{theorem}[Arity-polymorphic search specification]\label{thm:searchk}
For every fixed \(k\) and body \(f\), iPRA proves
\[
 \mu_f(N,\vec x)\le N,
\]
\[
 \mu_f(N,\vec x)<N\;\to\; f(\mu_f(N,\vec x),\vec x)=0,
\]
and
\[
 i<\mu_f(N,\vec x)\;\to\;f(i,\vec x)\ne0.
\]
\end{theorem}

\begin{proof}
Use quantifier-free induction on \(N\) with the conjunction of the three displayed statements as the
invariant; \(i\) and \(\vec x\) are parameters.  The zero case is immediate.  For the successor
case, split on \(\mu_N<N\) or \(\mu_N=N\).  In the first branch the defining equation gives
\(\mu_{N+1}=\mu_N\), so all clauses follow from the induction hypothesis.  In the second branch,
split on \(f(N,\vec x)=0\) or nonzero.  In the zero branch, \(\mu_{N+1}=N\); in the nonzero branch,
\(\mu_{N+1}=N+1\).  Order, equality-bit, and sign lemmas already used by the one-parameter proof
establish each subcase.  The formula is quantifier free because all functions are represented by
terms and all connectives are propositional.
\end{proof}

The Lean proof can be obtained by refactoring the current \code{Search.lean}: replace the two
specialized application helpers by \code{app} with a \code{Fin.cases} vector, and replace the
hard-coded simplification list by lemmas for the generic renaming \(\rho\).  The mathematical
induction is identical.

\subsection{Universal and existential aggregators at arbitrary arity}

Define a product recursion
\[
 \mathsf{allB}_f(N,\vec x)=\prod_{i<N}\operatorname{sg}(f(i,\vec x))
\]
and
\[
 \mathsf{exB}_f(N,\vec x)=1\mathbin{\dot-}
   \mathsf{allB}_{1\mathbin{\dot-}f}(N,\vec x).
\]
Here \(\dot-\) is truncated subtraction.  The proof that the universal bit equals the equality test
\([\mu_f(N,\vec x)=N]\) is again quantifier-free induction.  The existing one-parameter proof then
generalizes verbatim.

\begin{theorem}[Arity-polymorphic bounded reflection]\label{thm:allk}
For a fixed body \(f(i,\vec x)\), iPRA proves:
\begin{align*}
 \mathsf{allB}_f(N,\vec x)=1
 &\quad\Longleftrightarrow\quad
 \forall i<N\;f(i,\vec x)\ne0,\\
 \mathsf{exB}_f(N,\vec x)=1
 &\quad\Longleftrightarrow\quad
 \exists i<N\;f(i,\vec x)\ne0,
\end{align*}
where each direction is represented by an ordinary object derivation.  In the existential case, the
least witness is primitive recursive and satisfies the expected specification.
\end{theorem}

\begin{proof}
The equality \(\mathsf{allB}= [\mu=N]\) reduces the universal statement to the search
specification.  If the bit is one, then \(\mu=N\), so every \(i<N\) is below \(\mu\) and hence is not
a zero.  Conversely, if the bit were zero, the search specification would produce a zero below
\(N\), contradicting the universal premise.  The existential bit is the complement of the universal
bit for the complement body.  Its least witness is the least zero of that complement, which is
exactly the least nonzero of \(f\).
\end{proof}

This extension resolves the recurring ``one packed environment'' inconvenience while retaining the
existing API as the case \(k=1\).

\subsection{A reflected-formula interface}

For theorem generation, the useful abstraction is not a universal formula code but a package for one
fixed formula.

\begin{definition}[Reflected formula]
A reflected formula over \(n\) free variables is a tuple
\[
  (\varphi,b,C,S,T)
\]
where \(\varphi:\Fm\,\PTm\,n\), \(b:\PTm n\), and the object proofs establish:
\begin{enumerate}
\item \(C:\;b=0\lor b=1\);
\item \(S:\;b=1\to\varphi\) (soundness);
\item \(T:\;\varphi\to b=1\) (completeness).
\end{enumerate}
\end{definition}

The package is produced by meta-level recursion on a fixed bounded formula.  It contains the concrete
bit term and concrete object derivations; no formula code is evaluated inside iPRA.

\subsection{Propositional combinators}

Suppose \((\varphi,a)\) and \((\psi,b)\) are reflected.

\paragraph{Conjunction.}
Use \(a\cdot b\).  By case analysis on the two bit dichotomies, the product is one exactly when both
bits are one.  Soundness applies the two component soundness proofs; completeness applies the two
component completeness proofs and multiplication by one.

\paragraph{Disjunction.}
Use \(\operatorname{sg}(a+b)\).  If either formula holds, the corresponding bit is one, so the sum is
nonzero and its sign is one.  Conversely, case analysis on \(a,b\in\{0,1\}\) shows that a nonzero sum
forces at least one bit to be one; object-level disjunction elimination then yields
\(\varphi\lor\psi\).

\paragraph{Negation.}
Because \(\varphi\) is decidable through its bit, use \(1\mathbin{\dot-}a\).  If the complement bit is
one, then \(a=0\); assuming \(\varphi\) would give \(a=1\), contradiction.  Conversely, if
\(\neg\varphi\), the case \(a=1\) contradicts soundness, hence \(a=0\), and the complement is one.
This is intuitionistically valid because the reflected package supplies decidability.

\paragraph{Implication.}
For decidable antecedent \(\varphi\), \(\varphi\to\psi\) is equivalent intuitionistically to
\(\neg\varphi\lor\psi\).  Compile it using the previous two constructions.  No classical law is
used beyond the explicit decidability proof carried by the bit.

\subsection{Bounded quantifier combinators}

Let \(\theta(i,\vec x)\) have a reflected bit \(f(i,\vec x)\).  Apply
\cref{thm:allk} to obtain a bit for
\(\forall i<N\,\theta(i,\vec x)\).  Soundness first reflects the universal bit to
\(f(i,\vec x)\ne0\) at an arbitrary \(i<N\); since \(f\) is a bit, nonzero implies one, and the body
soundness proof gives \(\theta\).  Completeness starts from the bounded universal formula, applies
the body completeness theorem at arbitrary \(i\), and invokes universal-bit introduction.

The existential case uses \(\mathsf{exB}\).  Soundness extracts the least witness and applies body
soundness.  Completeness opens an object-level existential witness, applies body completeness, and
uses existential-bit introduction.  If the witness must not escape, the resulting existential is
sealed.  If it must escape, the least-witness component is retained.

\begin{theorem}[Static bounded-formula compilation]\label{thm:bounded-compile}
Every fixed bounded formula built from primitive-recursive atomic equations, conjunction,
disjunction, implication, negation, and bounded quantifiers admits a primitive-recursive bit and
object derivations of soundness, completeness, and bit-valuedness.
\end{theorem}

\begin{proof}
Meta-level structural induction on the fixed formula.  Atomic equality uses the checked equality bit;
atomic nonzero uses the sign bit.  The propositional cases are the constructions above.  The bounded
quantifier cases use the arity-polymorphic reflection theorem.  Every induction step constructs a
finite \(\Prf\) term.  This is a theorem schema implemented by Lean recursion, not an object-level
universal truth predicate.
\end{proof}

\subsection{Why this respects the limits of iPRA}

The compilation theorem does not say that arbitrary encoded formulas have a primitive-recursive
truth function.  For each fixed source formula, it constructs a specialized bit and proof.  Formula
size and quantifier nesting are consumed at elaboration time.  The object theory sees only a finite
term assembled from fixed primitive-recursive symbols.

Likewise, boundedness is essential.  An unbounded existential cannot in general be converted to a
primitive-recursive least witness.  It may still be eliminated locally by \cref{thm:seal}, but public
naming requires a proved cutoff.
